%!TEX program = xelatex
% ============================================================
%  הרצאה 2 — מבוא לכוחות וחוקי ניוטון
% ============================================================
\documentclass[12pt, a4paper]{article}

\usepackage{amsmath, amssymb, mathtools}
\usepackage{xcolor}
\usepackage{graphicx}

\usepackage{tcolorbox}
\tcbuselibrary{skins, breakable}
\tcbset{
  resultbox/.style={
    enhanced, colback=white, colframe=black,
    boxrule=0.8pt, arc=3pt,
    left=2pt, right=2pt, top=1.5pt, bottom=1.5pt,
  },
  notebox/.style={
    enhanced, breakable,
    colback=gray!8, colframe=gray!50,
    boxrule=0.8pt, arc=3pt,
    left=2pt, right=2pt, top=1.5pt, bottom=1.5pt,
    fonttitle=\bfseries,
  },
  warnbox/.style={
    enhanced, breakable,
    colback=red!5, colframe=red!60,
    boxrule=0.8pt, arc=3pt,
    left=3pt, right=3pt, top=2pt, bottom=2pt,
    fonttitle=\bfseries,
  },
  defbox/.style={
    enhanced, breakable,
    colback=blue!5, colframe=blue!50,
    boxrule=0.8pt, arc=3pt,
    left=3pt, right=3pt, top=2pt, bottom=2pt,
    fonttitle=\bfseries,
  },
}

\usepackage{tikz}
\usetikzlibrary{arrows.meta, calc, decorations.pathreplacing}
\usepackage[a4paper, top=2cm, bottom=2cm,
            left=2cm, right=2cm]{geometry}
\usepackage{setspace}
\setstretch{1.3}
\usepackage{titlesec}
\titleformat{\section}{\large\bfseries}{\thesection.}{0.5em}{}[\titlerule]
\titleformat{\subsection}{\bfseries}{\thesubsection}{0.5em}{}

\setlength{\headheight}{15pt}
\usepackage{fancyhdr}
\pagestyle{fancy}
\fancyhf{}
\rhead{\textenglish{\textbf{Lecture 2 --- Newton's Laws}}}
\lhead{\thepage}
\renewcommand{\headrulewidth}{0.3pt}

\usepackage{fontspec}
\usepackage{polyglossia}
\setmainlanguage{hebrew}
\setotherlanguage{english}
\setmainfont{Times New Roman}
\newfontfamily\hebrewfont{Times New Roman}
\newfontfamily\englishfont{Times New Roman}
\newfontfamily\hebrewfonttt{Courier New}

\newcommand{\hvec}[1]{\overrightarrow{#1}}

% ============================================================
\begin{document}
% ============================================================

% ── Title ────────────────────────────────────────────────────
\begin{center}
  {\LARGE\bfseries הרצאה 2 --- מבוא לכוחות וחוקי ניוטון}\\[6pt]
  \rule{\linewidth}{0.8pt}
\end{center}
\vspace{0.5cm}

% ============================================================
\section*{1. מבוא: מהו כוח?}
% ============================================================

הקינמטיקה (הרצאה 1) תיארה תנועה אך לא הסבירה \textbf{מה גורם לה}. זהו תפקיד הדינמיקה.

\begin{tcolorbox}[warnbox, title={בעיית ההגדרה המעגלית}]
ניסיון נאיבי: \emph{"כוח הוא מה שגורם לתאוצה"}.\\
אבל אז התאוצה מוגדרת על-ידי כוח, וכוח על-ידי תאוצה --- \textbf{הגדרה מעגלית}.\\[3pt]
פתרונו של ניוטון: לקבוע \textbf{חוקים} (אקסיומות) המאפיינים את התנהגות הכוחות, ולאמת אותם אמפירית.
\end{tcolorbox}

% ============================================================
\section*{2. החוק השלישי של ניוטון}
% ============================================================

\begin{tcolorbox}[defbox, title={חוק שלישי}]
אם גוף \textenglish{$A$} מפעיל כוח \textenglish{$\hvec{F}$} על גוף \textenglish{$B$},\\
אזי גוף \textenglish{$B$} מפעיל כוח \textenglish{$-\hvec{F}$} על גוף \textenglish{$A$}.\\[3pt]
שני הכוחות שווים בגודלם, מנוגדים בכיוונם, ופועלים על \textbf{גופים שונים}.
\end{tcolorbox}

\subsection*{2.1 דוגמה: קופסה על שולחן}

על הקופסה פועלים שני כוחות:
\textenglish{$m\hvec{g}$} כלפי מטה (משיכת כדור הארץ), ו-\textenglish{$\hvec{N}$} כלפי מעלה (כוח הנורמל מהשולחן).

\begin{center}
\begin{tikzpicture}[scale=1.3, >={Stealth[length=2.5mm]}]
  % surface
  \draw[thick] (-2.5, 0) -- (2.5, 0);
  \foreach \x in {-2.3,-1.9,-1.5,-1.1,-0.7,-0.3,0.1,0.5,0.9,1.3,1.7,2.1} {
    \draw[thick] (\x, 0) -- (\x - 0.25, -0.3);
  }
  % block
  \draw[thick, fill=blue!15] (-0.7, 0) rectangle (0.7, 1.2);
  \fill (0, 0.6) circle (1.6pt);
  % N arrow (up)
  \draw[->, very thick, green!55!black] (0, 0.6) -- (0, 2.1);
  \node[green!55!black, right] at (0.08, 1.7) {\textenglish{$\hvec{N}$}};
  % mg arrow (down)
  \draw[->, very thick, red] (0, 0.6) -- (0, -0.9);
  \node[red, right] at (0.08, -0.45) {\textenglish{$m\hvec{g}$}};
\end{tikzpicture}
\end{center}

\begin{tcolorbox}[warnbox, title={שגיאה נפוצה}]
\textenglish{$\hvec{N}$} ו-\textenglish{$m\hvec{g}$} שווים בגודלם ומנוגדים, אבל \textbf{הם אינם זוג של חוק שלישי}!\\[3pt]
שניהם פועלים על \textbf{אותו גוף} (הקופסה). חוק שלישי מקשר תמיד בין כוחות הפועלים על \textbf{שני גופים שונים}.
\end{tcolorbox}

\subsection*{2.2 הזוג הנכון של חוק שלישי}

הכוח \textenglish{$m\hvec{g}$} שכדור הארץ מפעיל על הקופסה מותאם, לפי החוק השלישי, לכוח שווה ומנוגד שהקופסה מפעילה על כדור הארץ.

\begin{center}
\begin{tikzpicture}[scale=1.2, >={Stealth[length=2.5mm]}]
  % Earth
  \draw[thick, fill=blue!25] (0, 0) circle (1.2);
  \node at (0, -0.05) {\small כדור הארץ};
  % Ball above
  \fill[gray!70!black] (0, 3.2) circle (0.18);
  \node[right] at (0.25, 3.2) {קופסה};
  % Force on ball (downward, from ball toward Earth)
  \draw[->, very thick, red] (0, 3) -- (0, 1.9);
  \node[red, right] at (0.1, 2.45) {\textenglish{$m\hvec{g}$}};
  % Force on Earth (upward, from Earth's center toward ball)
  \draw[->, very thick, blue!70!black] (0, 0) -- (0, 1.0);
  \node[blue!70!black, right] at (0.1, 0.5) {\textenglish{$-m\hvec{g}$}};
\end{tikzpicture}
\end{center}

(ההשפעה על תנועת כדור הארץ זניחה, בגלל המסה האדירה שלו.)

% ============================================================
\section*{3. מתקף הכוח (\textenglish{Impulse})}
% ============================================================

\begin{tcolorbox}[defbox, title={מתקף}]
המתקף של כוח \textenglish{$\hvec{F}(t)$} בקטע זמן \textenglish{$[t_1, t_2]$} הוא האינטגרל:
\[
  \hvec{I} = \int_{t_1}^{t_2} \hvec{F}(t)\,dt
\]
\end{tcolorbox}

המתקף הוא \textbf{וקטור}, ויחידותיו \textenglish{N$\cdot$s}.

% ============================================================
\section*{4. מערכת סגורה: כוחות פנימיים וחיצוניים}
% ============================================================

\begin{tcolorbox}[defbox, title={מערכת סגורה}]
אוסף גופים המשפיעים זה על זה בלבד --- ללא כוחות חיצוניים.
\end{tcolorbox}

\textbf{הבחנה בין סוגי כוחות:}
\begin{itemize}
  \item \textbf{כוח פנימי} --- בין שני גופים השייכים למערכת.
  \item \textbf{כוח חיצוני} --- מופעל ממקור מחוץ למערכת.
\end{itemize}

\medskip

\textbf{החוק השלישי בניסוח של מתקפים:} עבור זוג גופים \textenglish{$A, B$} במערכת סגורה,

\begin{tcolorbox}[resultbox]
\[
  \hvec{I}_{A \to B}(t) = -\hvec{I}_{B \to A}(t)
  \qquad \text{בכל זמן}
\]
\end{tcolorbox}

המתקף שמפעיל \textenglish{$A$} על \textenglish{$B$} שווה ומנוגד למתקף ש-\textenglish{$B$} מפעיל על \textenglish{$A$}.

% ============================================================
\section*{5. החוק השני של ניוטון}
% ============================================================

מתצפיות ניסיוניות: הכוח הכולל הפועל על גוף יוצר תאוצה \textbf{פרופורציונית} לו, באותו כיוון:
\[
  \hvec{F} \;\propto\; \hvec{a}
\]

מקדם הפרופורציה הוא תכונה של הגוף וקרוי \textbf{מסה}.

\begin{tcolorbox}[resultbox, title={החוק השני}]
\[
  \hvec{F} = m\hvec{a}
\]
\end{tcolorbox}

\textbf{משמעות הפרופורציה:} באותו כוח, גוף עם מסה גדולה יותר יתאיץ \textbf{פחות} --- \textenglish{$\hvec{a} = \hvec{F}/m$}.

% ============================================================
\section*{6. מסה}
% ============================================================

\begin{tcolorbox}[defbox, title={מסה}]
\textbf{תכונה סקלרית חיובית} של גוף, המבטאת את התנגדותו לשינוי במצב תנועתו (\textbf{אינרציה}).\\[3pt]
ככל שהמסה גדולה יותר, נדרש כוח גדול יותר כדי לייצר אותה תאוצה.
\end{tcolorbox}

\textbf{תכונות:}
\begin{itemize}
  \item \textbf{אדיטיביות (בקירוב טוב):} מסת מערכת = סכום מסות חלקיה. הקירוב נובע מהאטומיות של החומר.
  \item המסה היא תכונה אינהרנטית של הגוף ואינה תלויה במיקומו או בתנועתו.
\end{itemize}

% ============================================================
\section*{7. יחידות}
% ============================================================

\textbf{יחידת מסה:} ק"ג (\textenglish{kg}) במערכת \textenglish{SI}.

\textbf{יחידת כוח:} מהחוק השני,
\[
  [F] = [m]\cdot[a] = \text{kg}\cdot\text{m}\cdot\text{s}^{-2}
\]

\begin{tcolorbox}[resultbox]
\[
  1\;\text{N} \equiv 1\;\text{kg}\cdot\text{m}\cdot\text{s}^{-2}
  \qquad (\text{ניוטון})
\]
\end{tcolorbox}

\end{document}